% VisionDocs @ WACV 2027 -- SHORT PAPER DRAFT (4 pages + references), double-blind review copy.
% DRAFT ONLY. Not submitted, uploaded or posted. Draft of 22 Oct 2026.
%
% Template: official WACV 2027 Author Kit (wacv.sty, ieeenat_fullname.bst), downloaded unchanged from
%   https://wacv.thecvf.com/static/wacv-2027-author-kit-template.zip (linked from
%   https://wacv.thecvf.com/Conferences/2027/AuthorGuides). VisionDocs asks for "the standard main
%   conference WACV 2027 template". The kit's preamble.tex is inlined below.
% Build: pdflatex main; bibtex main; pdflatex main; pdflatex main  (not compiled here: no TeX engine
%   on this machine; see README.md).
% Numbers: every number has a "% src:" comment; README.md maps each claim to its source file.
%   Strip all comments before uploading LaTeX source anywhere (they name internal files).
% Anonymity: no names, affiliations, repository URLs or project links. Keep it that way.

\documentclass[10pt,twocolumn,letterpaper]{article}

%%%%%%%%% PAPER TYPE
% The kit requires a track option in review mode. Workshops have no tracks; "applications" is the
% closest fit. TODO(owner): confirm with the organisers or keep.
\usepackage[review,applications]{wacv}      % REVIEW version
% \usepackage{wacv}                         % CAMERA-READY version
% \usepackage[pagenumbers]{wacv}            % arXiv version with page numbers

% ---- inlined from the kit's preamble.tex
\newcommand{\red}[1]{{\color{red}#1}}
\newcommand{\TODO}[1]{\textbf{\color{red}[TODO: #1]}}
% \renewcommand{\TODO}[1]{}   % uncomment to hide TODOs
% ---- end of preamble

\definecolor{wacvblue}{rgb}{0.21,0.49,0.74}
\usepackage[pagebackref,breaklinks,colorlinks,allcolors=wacvblue]{hyperref}

\def\wacvPaperID{*****} % TODO(owner): VisionDocs short papers go through a Google Form, which may not assign an ID
\def\confName{WACV}
\def\confYear{2027}

\newcommand{\sysname}{Section Anything\xspace} % one place to rename the system for review if needed
\newcommand{\na}{--}

\title{\sysname: Heading Hierarchies for Structure-Aware Retrieval\\from the Text Layer of Born-Digital PDFs}

\author{Anonymous}  % review mode prints "Anonymous WACV ... submission" instead; never add names here

\begin{document}
\maketitle

%%%%%%%%%%%%%%%%%%%%%%%%%%%%%%%%%%%%%%%%%%%%%%%%%%%%%%%%%%%%%%%%%%%%%%%%%%%%%%%%%%%%%%%%%%%%%%%%%%%%
\begin{abstract}
Retrieval pipelines increasingly split documents at their headings, so a parser that returns flat
or wrong heading levels changes what every chunk contains and which section path it carries. We
study heading-hierarchy recovery for born-digital PDFs from the text layer alone: typography,
numbering, bookmarks and printed contents, optionally reviewed by a small language model whose
proposals must anchor to verbatim text. The method performs no OCR and uses no page images. On 905
READoc arXiv papers not used during development, scored with READoc's own Heading Tree TEDS, the
full-text setting scores 70.0 against 49.4 for Docling with its heading-hierarchy stage enabled and
64.8 for MinerU with its own LLM title levelling. On GitHub READMEs it only ties Docling. The model settings cost
US\$0.07--0.16 per 1,000 pages with a small model. As a second contribution, an opt-in bookmark check
for Docling's heading stage, which demotes section headers that a complete PDF outline does not
list, adds 5.4 TEDS points on 100 arXiv papers and 3.7 on 100 READMEs, and does nothing on a
pre-registered set of legislation. On Qasper, section-aware chunks retrieved the evidence more often,
but answer quality did not change.
\end{abstract}
% src: 905/70.0/49.4 -- site/data/results.js readoc.heldout; HANDOFF "Final held-out result, arXiv (21 Oct)". MinerU 64.8 -- results.js readoc.heldoutMinerU; .tmp/readoc-pilot/heldout_report_mineru_full.txt (25 Oct).
% src: GitHub tie -- results.js readoc.full.note. Cost -- results.js modes[].costPer1000Pages.
% src: Docling option -- docs/docling-pr/results.md "Main result" and "Fresh documents (22 Oct)". Qasper -- docs/rag-qasper.md.

%%%%%%%%%%%%%%%%%%%%%%%%%%%%%%%%%%%%%%%%%%%%%%%%%%%%%%%%%%%%%%%%%%%%%%%%%%%%%%%%%%%%%%%%%%%%%%%%%%%%
\section{Introduction}
\label{sec:intro}

Retrieval-augmented generation (RAG) over long documents depends on chunking. A growing number of
pipelines cut along the document's own sections, or prefix each chunk with its section path, so that
a retrieved passage keeps its context~\cite{lu2026hichunk,anthropic2024contextual}. LightRAG's
heading-based chunker, for example, splits on the heading levels a parser supplies~\cite{guo2025lightrag,lightrag_parser_docs},
and Markdown header splitters do the same with \texttt{\#} levels. The recovered heading
hierarchy then decides the chunks: a missed heading merges two topics, a sentence promoted to a
heading splits a paragraph, and a wrong level files a subsection under the wrong parent.
% TODO(owner): re-verify the LightRAG chunker wording against its current docs (refs.bib entry is marked UNVERIFIED).

Open-source converters such as Docling~\cite{auer2024docling,docling_software} and
MinerU~\cite{wang2024mineru,mineru_software} detect section-header blocks on page images and then
assign depth with rules or an optional LLM step. In the READoc benchmark~\cite{li2025readoc}, the
published Heading Tree TEDS of Docling, MinerU and Marker on arXiv papers ranged from 38.3 to 42.0,
while Nougat~\cite{blecher2024nougat}, an OCR model trained on arXiv, reached 88.5. Born-digital PDFs
also carry signals that a page image shows only indirectly or not at all: an outline (bookmarks),
font names and exact sizes, and the text itself with its numbering. None is reliable alone: bookmarks
are often missing or partial, and bold type also marks theorem labels, FAQ questions and table cells.
% src: published READoc scores -- results.js readoc.published (Docling 38.3, Marker 39.4, MinerU 42.0, Nougat-base 88.5); checked against READoc Tables 3-4 per paper.md section 1.

We ask how far these text-layer signals go, and whether they can repair a vision pipeline.
Our contributions are:
(i)~a heading-hierarchy method for born-digital PDFs with three cost settings, in which every
model proposal must anchor to a verbatim span (\cref{sec:method});
(ii)~an evaluation with READoc's own scorer on held-out arXiv papers and on GitHub READMEs, against
current Docling and MinerU, including MinerU's own LLM title levelling (\cref{sec:readoc});
(iii)~an opt-in bookmark check for Docling, tested on real conversions, on partial outlines and on
documents not used in its design, one set pre-registered (\cref{sec:docling}); and
(iv)~a Qasper RAG test in which retrieval improved but answers did not (\cref{sec:rag}).
The scope is narrow on purpose: everything here needs a PDF text layer. We perform no OCR, and
scanned pages yield nothing.

\paragraph{Related work.} Besides READoc, HRDoc~\cite{ma2023hrdoc}, Comp-HRDoc~\cite{wang2024detectorderconstruct}
and DocHieNet~\cite{xing2024dochienet} evaluate hierarchical structure, mainly for layout and vision
models; we have not yet evaluated on them. PyMuPDF4LLM maps larger font sizes to Markdown header
levels or reads the PDF's table of contents~\cite{pymupdf4llm_software}. Docling's optional
heading-hierarchy stage sets header levels from bookmarks, numbering and font style, and is off by
default~\cite{docling_software}.

%%%%%%%%%%%%%%%%%%%%%%%%%%%%%%%%%%%%%%%%%%%%%%%%%%%%%%%%%%%%%%%%%%%%%%%%%%%%%%%%%%%%%%%%%%%%%%%%%%%%
\section{Method}
\label{sec:method}

\paragraph{Text layer.} We read lines in reading order with pdfminer's layout analysis. Each line
keeps its text, box, median glyph size and a bold flag. Boldness comes from font names, including
the codes of \TeX{} bold faces (CMBX, SFBX, \etc); without them, the bold headings of almost every
\LaTeX{} paper looked like body text. Bookmarks are read with their level, title, target page and
target height. Shared detectors mark printed tables of contents and page furniture (running headers,
page numbers), which may never start a section. Extraction takes about 0.02\,s per page.
% src: paper.md 3.1-3.2; extract.py; speed HANDOFF L148-150 (via paper.md 3.1).

\paragraph{Free (no model).} The outline, or else the printed table of contents, is anchored to body
lines by title on its target page; it becomes the backbone when at least three entries and 60\% of
the outline anchor. Typography and numbering rules add the remaining headings, after filters for
captions, theorem labels, bold mathematics, repeated callout labels and bold run-in terms. Levels come
from dotted numbering depth (``2.3.1'' is level 3) or else from style rank.
% src: paper.md 3.3; Appendix C.

\paragraph{Quick check.} A looser filter proposes candidate lines. A small model sees each as one
row with size and bold tags and the first 70 characters of the next line, in batches of 300 with the
last eight decisions as context, and returns heading or not, level, number and title through a
schema-constrained call. A title not found in print is replaced by the printed line. Quick check
cannot recover a heading that no rule proposed.
% src: paper.md 3.4.

\paragraph{Standard.} The model reads every line, numbered and tagged, in windows of about 120,000
characters that end at page breaks, with the last eight sections as context. For each section it
returns the start line, number, title, level and a \emph{start marker}: the first 8--15 words,
copied exactly (and an end marker when the section ends inside the window).
% src: paper.md 3.5.

\paragraph{Resolution.} Models rarely quote byte for byte, so both the reply and the document are
reduced to a canonical form (case-folded, ligatures expanded, line-break hyphens removed) with a map
back to original offsets. Markers are matched near the model's line hint, then anywhere; matches
inside a table of contents are dropped when a body copy exists; a bounded fuzzy fallback is flagged.
Proposals with the fewest candidate positions claim offsets first, so no two sections share a start.
End markers and printed titles are checked, and mismatches are flagged rather than accepted. Each
section thus has a level, number, title, exact character range, pages and line boxes. Finally, the
first heading is flagged as the document title if it is on page~1, unnumbered, not a generic name
and at least 1\,pt larger than every other heading. Prompts treat document text as untrusted, and a
per-run cost cap is enforced.
% src: paper.md 3.6-3.8.

%%%%%%%%%%%%%%%%%%%%%%%%%%%%%%%%%%%%%%%%%%%%%%%%%%%%%%%%%%%%%%%%%%%%%%%%%%%%%%%%%%%%%%%%%%%%%%%%%%%%
\section{Heading trees on READoc}
\label{sec:readoc}

\paragraph{Setup.} We use READoc's arXiv subset (1,003 papers whose keys have headings) and its
GitHub subset (1,224 READMEs rendered to PDF); the Zenodo subset has no published keys. Each system's
Markdown goes through READoc's own cleaning and segmentation code and is scored with its Heading Tree
TEDS ($\times$100). Docling~2.140 runs with its heading-hierarchy stage enabled and is scored on its
own Markdown export. MinerU runs on its cloud service (3.4.4, VLM mode); for ``+\,LLM levels'' we
applied MinerU's own LLM title-levelling step to that output with gpt-6-luna (Luna) or gpt-6.1-sol
(Sol). Our model settings use Luna at low reasoning effort.
% src: paper.md 4.1-4.2; HANDOFF L832, L864, L1079; MinerU relay HANDOFF L607-621.

\paragraph{Protocol.} A 200-document sample (100 per subset, fixed seed) guided every method change,
so our main test is the 905 \emph{held-out} arXiv papers with headings. MinerU's cloud service, on its
free quota, processed all of them, so every system is scored on the same papers. Paired comparisons use per-document
differences with a 95\% bootstrap interval (4,000 resamples) and a Wilcoxon signed-rank test. Each
configuration was run once.
% src: paper.md 4.3; heldout_report.py (via paper.md).

\begin{table}[t]
\centering
\footnotesize
\setlength{\tabcolsep}{4pt}
\begin{tabular}{@{}lcc@{}}
\toprule
 & arXiv, held out & GitHub \\
System & all 905 & all 1,224 \\
\midrule
Ours: Standard (Luna)        & \textbf{70.0} & \textbf{75.2} \\
Ours: Quick check (Luna)     & 63.9 & 74.3 \\
Ours: Free (no model)        & 52.0 & 75.0 \\
Docling 2.140, hierarchy on  & 49.4 & 75.0 \\
MinerU + LLM levels (Sol)    & 64.8 & \na \\
MinerU + LLM levels (Luna)   & 62.7 & \na \\
MinerU cloud (3.4.4, VLM)    & 53.9 & \na \\
\bottomrule
\end{tabular}
\caption{READoc Heading Tree TEDS ($\uparrow$). arXiv: papers not used during development, all
processed by MinerU's cloud. GitHub includes the 100 development READMEs; on the
1,124 held-out READMEs all our settings tie Docling (see text). Luna = gpt-6-luna, Sol = gpt-6.1-sol.}
\label{tab:readoc}
\end{table}
% src: col 905 -- results.js readoc.heldout.note (70.0/63.9/49.4); Free 52.0 -- HANDOFF "Final held-out result" (all 905).
% src: MinerU rows -- results.js readoc.heldoutMinerU; .tmp/readoc-pilot/heldout_report_mineru_full.txt (25 Oct).
% src: col GitHub -- results.js readoc.full (latest, 21 Oct); HANDOFF "Full READoc, latest version (21 Oct)".

\paragraph{Results.} On all 905 held-out papers, Standard leads Docling by 20.6 points (95\% CI 19.3
to 21.8), better on 773 papers and worse on 103 (\cref{tab:readoc}). On the same papers it leads
MinerU with LLM levels on the larger model by 5.1 points (95\% CI 4.0 to 6.2; $p=1.3\times10^{-21}$)
and on the same small model by 7.2 (6.2 to 8.2). MinerU's own LLM step matters: it lifts MinerU from
53.9 to 64.8, although with Sol it gave up on 154 papers (no parsable reply in three tries) and kept
the raw levels; without them our lead is 2.8 (1.6 to 3.9).
Quick check, which sends only candidate lines, falls between MinerU's two LLM variants. A prompt fix
for papers (the title, Abstract and References are headings; theorem labels never are) was designed
on the development sample, where it added 5.7 points to Standard; on the held-out papers it added
5.1, so it appears to generalise.
% src: HANDOFF "Final held-out result" (+20.6, CI, 773/103, +5.1 prompt fix); MinerU (+5.1, CI, p, +7.2, 154, +2.8) -- .tmp/readoc-pilot/heldout_report_mineru_full.txt (25 Oct), results.js readoc.heldoutMinerU.note; paper.md 4.4 (+5.7 on sample).

\paragraph{Where we do not win.} On the 1,124 held-out READMEs, every setting ties Docling: Standard
$+0.3$ (95\% CI $-0.9$ to $+1.6$), Quick check $-0.6$ ($-1.8$ to $+0.7$), Free $+0.1$ ($-1.0$ to
$+1.2$). On the 100 development READMEs, Docling (73.2) scored above Standard (71.3), while MinerU
(48.2 to 56.5) and PyMuPDF4LLM (57.5) trailed. Free's small arXiv lead comes from bookmarks: on
papers without them it scores 42.3 against Docling's 46.7. Nougat's published 88.5 on arXiv remains
far above every number here.
% src: held-out GitHub deltas -- paper.md 4.6 (.tmp/readoc-github-latest/scores.json) and HANDOFF "Full READoc, latest version".
% src: dev README scores -- results.js readoc.sample. Bookmark split -- HANDOFF L801-802. Nougat -- results.js readoc.published.

\begin{table}[t]
\centering
\footnotesize
\setlength{\tabcolsep}{4pt}
\begin{tabular}{@{}lccc@{}}
\toprule
Setting & US\$ / 1,000 pp & F1 & Levels \\
\midrule
Ours: Free                     & 0    & 0.800 & 0.977 \\
Ours: Quick check (Luna)       & 0.07 & 0.911 & 0.966 \\
Ours: Standard (Luna)          & 0.16 & \na   & \na \\
Ours: Standard (Sol)           & 2.9  & \textbf{0.920} & \textbf{0.990} \\
Docling 2.140, hierarchy on    & 0 (local) & 0.896 & 0.957 \\
MinerU cloud + LLM levels (Sol) & \na & 0.870 & 0.950 \\
\bottomrule
\end{tabular}
\caption{Cost per 1,000 pages, and heading F1 and level agreement on 23 public, non-scanned documents
of our multi-genre corpus (1,282 pages), from runs of 19 Oct 2026 that predate the final method. Model
settings also accept a local model at no cost; we have not measured its accuracy.}
\label{tab:cost}
\end{table}
% src: cost -- results.js modes[].costPer1000Pages (0, 0.07, 0.16) and costPer1000PagesStrong (2.9).
% src: F1/Levels (23 docs, no OpenStax) -- docs/PLAN-upstream-and-rag.md "Where we stand"; paper.md Table 4 (F1 23 / Levels 23); HANDOFF L674-677.

\paragraph{Other genres and cost.} To look beyond papers and READMEs we built an English corpus whose
answer keys come from source markup (legislation XML, \LaTeX{} source, Wikipedia's parse API,
generators) rather than from hand labelling. A predicted heading matches a key heading when number
and title are at least 85\% similar and it starts within one page; \emph{level agreement} is the share
of consecutive matched pairs whose relative depth agrees with the key. On the 23 public non-scanned
documents (\cref{tab:cost}), both model settings score above Docling and MinerU in F1 and level
agreement; this set is small and was run before the final method, so we treat it as a sanity check. Measured on
this set, Quick check cost \$0.04 and Standard (Sol) \$2.33 per 1,000 pages.
% src: corpus method -- paper.md 4.7; evaluate.py. Measured costs -- PLAN-upstream-and-rag.md "Where we stand".
% NOTE: three OpenStax documents are excluded from every corpus figure (licence forbids AI use). Do not add 26- or 32-document figures.

%%%%%%%%%%%%%%%%%%%%%%%%%%%%%%%%%%%%%%%%%%%%%%%%%%%%%%%%%%%%%%%%%%%%%%%%%%%%%%%%%%%%%%%%%%%%%%%%%%%%
\section{A bookmark check for a vision pipeline}
\label{sec:docling}

Docling's heading stage already uses the outline to set the levels of matched headers, but a section
header that matches no bookmark is always kept. Many such headers are not
headings (paper titles, author lines, algorithm captions, lead-ins such as ``Example:''). We added an
option, off by default, that demotes them to text when the outline looks complete: it has more than
one level, every dotted-number header (``2.1 \ldots'') matches a bookmark, at least half of the
headers match, and no two headers to be demoted share a heading-like style (larger, bold or a
non-body font face), since such a family usually means the outline skips a level. Standard headings
(Abstract, References, Appendix) are always kept and nothing is promoted. The layout model still
proposes the headers from the page image; the outline only vetoes the ones it does not list.
% src: docs/docling-pr/issue_draft.md "Requested feature"; results.md "Summary".

\begin{table*}[t]
\centering
\footnotesize
\setlength{\tabcolsep}{5pt}
\begin{tabular}{@{}llcccll@{}}
\toprule
Documents & Answer key / metric & $n$ & Off & On & $\Delta$ (95\% CI) & better / same / worse \\
\midrule
READoc arXiv, full outline           & READoc, TEDS & 100 & 50.4 & 55.8 & $+5.4$ ($+4.0$ to $+7.0$) & 42 / 56 / 2 \\
READoc GitHub, full outline          & READoc, TEDS & 100 & 75.7 & 79.4 & $+3.7$ ($+2.1$ to $+5.4$) & 27 / 67 / 6 \\
READoc arXiv, deepest level removed  & READoc, TEDS & 69  & 51.6 & 51.9 & $+0.3$ ($0.0$ to $+0.6$)  & 4 / 65 / 0 \\
READoc GitHub, deepest level removed & READoc, TEDS & 100 & 75.7 & 75.3 & $-0.4$ ($-1.0$ to $0.0$)  & 0 / 97 / 3 \\
READoc, top level only (both)        & READoc, TEDS & 200 & \na & \na & $0.0$                      & 0 / 200 / 0 \\
arXiv papers outside READoc          & \LaTeX{} source, TEDS & 8 & \na & \na & $+4.4$ ($+0.3$ to $+10.7$) & 4 / 4 / 0 \\
Canadian federal acts, pre-registered & official XML, F1 & 9 & \na & \na & $0.0$ (nothing demoted) & 0 / 9 / 0 \\
\bottomrule
\end{tabular}
\caption{Docling 2.140 with the bookmark check off and on, same conversion per document. Off is
identical to stock Docling on all 200 READoc documents. READoc rows use bookmarked documents outside
our development sample; the outline-removal rows test partial outlines.}
\label{tab:docling}
\end{table*}
% src: rows 1-5 -- docs/docling-pr/results.md "Main result" and "Risk check" (commit 97c09a7).
% src: rows 6-7 -- docs/docling-pr/results.md "Fresh documents (22 Oct)" (commits 97c09a7 then 183f1dd; scores unchanged by the fixes).
% "same" for the 8 papers = 8 - 4 better - 0 worse.

\paragraph{Results.} On real conversions of 200 READoc documents the check adds 5.4 points on arXiv
and 3.7 on GitHub (\cref{tab:docling}; Wilcoxon $p<10^{-4}$ for both). On arXiv, 57 of the 68
demoted headers are the paper title, which the layout model labels a section header; in 56 of those
papers READoc's key has that line as the document title. Relabelling a demoted first header as a
title instead of text raises the arXiv gain to $+8.4$ (95\% CI $+6.4$ to $+10.5$; GitHub $+3.4$).
With a partial outline the mean effect is small, but 3 of 100 READMEs lose 10 to 17 points: the
removed level held a single heading, which the style check cannot tell from a false header.
% src: docs/docling-pr/results.md "Summary", "What gets demoted", "Risk check", "Title variant".

\paragraph{Documents not used in the design.} The thresholds were chosen while looking at READoc, so we
also ran documents never used in the design. On 8 arXiv papers outside READoc (keys from their \LaTeX{}
source) the gain has the same sign ($+4.4$; $+5.5$ with title relabelling), none worse. On 9 Canadian
federal acts, pre-registered with keys, metrics and code frozen before any output, the layout model
labels marginal notes as section headers, far fewer than half of the headers match a bookmark, and
the check correctly does nothing. Two general fixes found on these fresh documents (letter-spaced
standard headings; bookmark titles that differ from the printed heading) left these scores unchanged;
the READoc rows predate them and have not been re-run, and only the Canadian acts were a blind test.
% src: docs/docling-pr/results.md "Fresh documents (22 Oct)".

%%%%%%%%%%%%%%%%%%%%%%%%%%%%%%%%%%%%%%%%%%%%%%%%%%%%%%%%%%%%%%%%%%%%%%%%%%%%%%%%%%%%%%%%%%%%%%%%%%%%
\section{Do better headings give better answers?}
\label{sec:rag}

We sampled 200 of the 415 Qasper~\cite{dasigi2021qasper} test papers with an answerable question
(fixed seed; 707 questions). Every arm used our text extraction, the same embeddings, answer model
(Luna) and context budget; only chunk boundaries and section-path prefixes changed: fixed
1,200-token chunks (A), Standard headings (C), the papers' gold section names (D) and Quick check
(E). Answers were scored with Qasper's official answer F1; evidence recall asks whether the retrieved
context contains a gold evidence paragraph (653 questions).
% src: docs/rag-qasper.md "Setup".

Section-aware chunks retrieved the evidence more often: 0.965 against 0.933 at a 6,000-token budget
($p=0.002$) and 0.649 against 0.496 at 1,200 tokens ($p<0.001$), with gold, Standard and Quick check
headings almost equal. A control with fixed chunks of the section chunks' average size recovered most
of that gain; section boundaries added a significant 8.7 points only at 1,200 tokens, and the
section-path prefix added nothing measurable. Answer F1 did not differ (A 0.541, C 0.534; C$-$A
$=-0.007$, 95\% CI $-0.030$ to $+0.015$), while re-asking identical excerpts moved F1 by 1.5 points.
Qasper papers are short (58 of 200 fit in 6,000 tokens), so we cannot claim that better headings give
better answers; long documents remain the open test. A 51-question FinanceBench~\cite{islam2023financebench}
pilot agrees: heading-split chunks found the evidence page more often than fixed chunks (34 against
27), but our headings did not beat MinerU's (30 against 31 correct answers).
% src: docs/rag-qasper.md "Evidence recall", "Answer F1", "Noise floor", "Setup"; paper.md 6.2.
% src: FinanceBench -- paper.md Table 9; results.js rag.rows; HANDOFF "RAG pilot on FinanceBench".

%%%%%%%%%%%%%%%%%%%%%%%%%%%%%%%%%%%%%%%%%%%%%%%%%%%%%%%%%%%%%%%%%%%%%%%%%%%%%%%%%%%%%%%%%%%%%%%%%%%%
\section{Limitations and conclusion}
\label{sec:limits}

\paragraph{Scope.} The method reads the PDF text layer. It performs no OCR, so scanned pages yield
nothing, and it does not look at page images, so headings marked only by colour or position are hard
to find. Our corpus is English only. On READMEs we only tie Docling, and without bookmarks the free
setting is weaker than Docling on papers.

\paragraph{Evidence.} Every configuration was run once with proprietary models (October 2026), and we
saw run-to-run variation on single documents. Method changes were designed on the READoc development
sample; held-out arXiv papers are the main test. The corpus results predate the final method. Marker
could not be run on our hardware, and PageIndex, HRDoc, Comp-HRDoc and DocHieNet are not yet covered.
The Docling check was tuned on READoc and tested only on papers, READMEs and legislation. Quick check
and Standard send text to the configured model provider; Free and local models do not. We excluded
textbooks whose licence forbids generative-AI use from all reported corpus results.
% src: paper.md 5.2, 5.3, 7; HANDOFF L674-677 (OpenStax).

\paragraph{Conclusion.} For born-digital PDFs, the text layer, its typography and the document's own
outline carry most of what a heading hierarchy needs; a small model, held to verbatim text, adds the
judgement that rules lack for cents per thousand pages. The same outline signal can veto false headers
in a vision pipeline. Whether better headings yield better RAG answers remains open. Code, scripts
and per-document scores will be made available after review.

{
    \small
    \bibliographystyle{ieeenat_fullname}
    \bibliography{refs}
}

\end{document}
